\section{规则三：期货截面统计套利}\label{sec:statarb}

\subsection{适用性依据}

论文表 5 的期货截面口径（XS）结果：

\begin{table}[htbp]
\centering\footnotesize
\caption{期货截面反转的实测（论文表 5，成本 2 基点）}\label{tab:stasrc}
\begin{tabular}{@{}>{\raggedright\arraybackslash}p{1.2cm}>{\raggedright\arraybackslash}p{1.7cm}>{\raggedright\arraybackslash}p{1.7cm}>{\raggedright\arraybackslash}p{1.7cm}>{\raggedright\arraybackslash}p{1.7cm}>{\raggedright\arraybackslash}p{1.7cm}>{\raggedright\arraybackslash}p{1.5cm}@{}}
\toprule
\textbf{$L$} & \textbf{毛利夏普} & \textbf{净夏普} & \textbf{年化收益} & \textbf{最大回撤} & \textbf{年化换手} & \textbf{NW $t$} \\
\midrule
5   & $-0.05$ & $-0.20$ & $-2.59\%$ & $-33.53\%$ & $79.58$ & $-0.76$ \\
10  & $-0.20$ & $-0.31$ & $-3.81\%$ & $-46.20\%$ & $57.81$ & $-1.19$ \\
20  & $+0.07$ & $-0.01$ & $-0.71\%$ & $-34.84\%$ & $41.76$ & $-0.04$ \\
60  & $+0.77$ & $+0.72$ & $+7.57\%$ & $-18.74\%$ & $26.63$ & $+2.59$ \\
120 & $+0.60$ & $+0.56$ & $+5.44\%$ & $-17.40\%$ & $20.78$ & $+1.98$ \\
250 & $+0.28$ & $+0.25$ & $+2.09\%$ & $-26.64\%$ & $16.08$ & $+0.82$ \\
\bottomrule
\end{tabular}
\end{table}

论文表 5 中的这组结果有一个值得单独强调的地位：$L=60$ 的 Newey--West $t$ 为 $2.59$，是论文全部 48 个组合中少数通过常规显著性门槛的结果。同时它的最大回撤仅 $-18.74\%$、成本占毛利约 $6.3\%$（在全部门槛检查中余量最充足）。

三点必须读出。

\textbf{其一，短窗口在这条口径上是负收益}。$L=5$ 与 $L=10$ 的毛利夏普已经是负的（$-0.05$、$-0.20$），$L=20$ 接近零（$+0.07$）。这与论文的 U 形结构判断一致：期货截面上的短线趋势（$L\le10$）与中线反转（$L\ge60$）之间存在一个过渡带，$L=20$ 恰在过渡带中央，两边都不占优。

\textbf{其二，$L=60$ 与 $L=120$ 构成可行区间}。二者的净夏普 $0.72$ 与 $0.56$、回撤 $-18.74\%$ 与 $-17.40\%$、$t$ 值 $2.59$ 与 $1.98$。$L=250$ 的收益与显著性都明显衰减。

\textbf{其三，本规则的实现方式与论文不同，须明确标注}。论文的截面口径是“按过去 $L$ 日累计收益排序，多后 20\%、空前 20\%”，这是一个\emph{纯截面排序}策略，不使用协整、不估计 OU、不做对对冲。而本文设计的规则三走向\emph{协整配对}路线，理由有两条：其一，纯截面排序在多空两端会引入大量方向性风险（组合对商品指数有净暴露），而配对天然对冲共同因子；其二，论文的对偶七指出反转策略在结构上应带“资格检验”，协整检验正是这道门槛，可以让规则在最弱的组合（$L=20$ 附近的过渡带）上直接弃权。

必须诚实说明：\textbf{配对版本的历史绩效不能直接引用论文表 5 的数字}——两者是不同的策略构造。论文表 5 支持的是“期货截面上中线反转方向有效”这一判断，本文据此选择 $L=60/120$ 作为信号窗口，但配对化的收益水平需自行回测，本文不给出预期数字。

\subsection{配对与篮子形成}

配对交易的方法论基础由 Vidyamurthy(2004) 的专著系统化，其覆盖的距离法、协整法与随机价差法三种配对选择思路至今仍是该领域的技术主干\cite{vidyamurthy2004}；Krauss(2017) 的综述梳理了该族策略的方法演进与利润衰减脉络\cite{krauss2017}。本规则的构造（产业链内配对 + 协整筛选 + OU 阈值）属于其中的协整分支。

\textbf{配对范围：产业链内}。跨产业链配对在商品期货上往往不成立（例如黑色与贵金属之间缺乏经济意义上的长期均衡关系）。建议的产业链划分：

\begin{table}[htbp]
\centering\footnotesize
\caption{建议的配对产业链划分}\label{tab:chains}
\begin{tabular}{@{}>{\raggedright\arraybackslash}p{3.2cm}>{\raggedright\arraybackslash}p{10.8cm}@{}}
\toprule
\textbf{产业链} & \textbf{候选品种} \\
\midrule
黑色系 & RB（螺纹）、HC（热卷）、I（铁矿）、J（焦炭）、JM（焦煤）、SF（硅铁）、SM（锰硅） \\
能源化工 & MA（甲醇）、PP（聚丙烯）、L（塑料）、TA（PTA）、EG（乙二醇）、V（PVC）、BU（沥青）、SC（原油）、FU（燃油） \\
有色金属 & CU（铜）、AL（铝）、ZN（锌）、NI（镍）、SN（锡）、PB（铅） \\
农产品 & M（豆粕）、Y（豆油）、P（棕榈）、C（玉米）、CS（淀粉）、OI（菜油）、RM（菜粕）、SR（白糖）、CF（棉花）、AP（苹果） \\
油脂油料 & A（豆一）、B（豆二）、M、Y、P、OI \\
\bottomrule
\end{tabular}
\end{table}

\textbf{形成期与交易期}。形成期 60 交易日，交易期 21 交易日，滚动推进。这一“形成期与交易期严格分离”的设计来自 Gatev 等(2006) 的 6+6 个月滚动范式\cite{gatev2006}，其目的是消除 look-ahead——若配对是在包含交易期的数据上选出的，价差平稳性会被机械地保证。

\textbf{候选对生成}。同一产业链内两两组合，形成期相关性 $\ge0.6$ 者进入候选池。

\subsection{协整资格检验（入场门槛）}

这道门槛是本规则区别于规则一、二的核心结构，其依据是论文对偶七：反转策略在结构上多一道“资格检验”，是“条件性”策略。

\textbf{两步法}（Engle--Granger\cite{englegranger1987}）：

\begin{enumerate}[leftmargin=2.2em]
  \item 对形成期数据做回归 $\ln P^A_t=\alpha+\beta\ln P^B_t+e_t$；
  \item 对残差 $e_t$ 做 ADF 检验\cite{dickeyfuller1979}，统计量需小于 5\% 临界值。
\end{enumerate}

不通过的对直接弃权，不进入交易。这道门槛的价值不在于提高单对收益，而在于\emph{减少无效样本}——论文指出反转策略的“资格检验”会消耗样本，使有效样本基础比趋势策略更薄，这是它的结构性代价。

\subsection{OU 参数估计与半衰期筛选}

对通过协整检验的价差 $e_t$ 拟合 AR(1)：

\begin{equation}
e_{t+1}=c+\phi e_t+u_t \;\Longrightarrow\;
\theta=-\frac{\ln\phi}{\Delta},\qquad
H=\frac{\ln 2}{\theta} .
\end{equation}

\textbf{半衰期筛选区间}：$H\in[5,\,60]$ 交易日。上界的含义是资金占用：$H$ 过大意味着价差需要数月才收敛一半，占用保证金的效率太低；下界的含义是成本：$H$ 过小意味着价差变化极快，交易频率上升、成本侵蚀收益。

\textbf{估计偏差的提醒}。论文附录指出 $\hat\phi$ 在小样本下有向下偏误，短样本估出的 $\theta$ 往往偏大、半衰期偏短\cite{kendall1953}。形成期 60 交易日对 AR(1) 估计而言偏短，因此建议对 $\hat\phi$ 做偏误修正，或把形成期延长至 120 交易日。这一点在本文中属\textbf{工程判断}。

\subsection{对冲比率}

\textbf{静态版本的缺陷}。若对冲比率 $\beta$ 在整个交易期保持不变（用形成期的 OLS 系数），则当两品种的相对关系发生漂移时，价差序列会失真——它会把参数漂移误读为价差偏离，产生虚假信号。

\textbf{在线版本（推荐）}。用 Kalman 滤波滚动估计 $\beta_t$，其状态空间形式为

\begin{equation}
\begin{aligned}
\ln P^A_t &= \alpha_t+\beta_t\ln P^B_t+\varepsilon_t, &&\varepsilon_t\sim\mathcal N(0,\sigma_\varepsilon^2),\\
\alpha_{t+1}&=\alpha_t+\eta^{\alpha}_t, \quad
\beta_{t+1}=\beta_t+\eta^{\beta}_t, &&\eta\sim\mathcal N(0,\Sigma_\eta),
\end{aligned}
\end{equation}

用时刻 $t$ 及之前的信息更新 $(\alpha_t,\beta_t)$，据此计算价差 $e_t=\ln P^A_t-\alpha_t-\beta_t\ln P^B_t$。这一做法来自 Elliott 等(2005) 的 OU 状态空间配对框架\cite{elliott2005}。

\subsection{交易带（本规则的核心）}

\textbf{标准化价差}：

\begin{equation}
z_t=\frac{e_t-\mu_t^{(60)}}{\sigma_t^{(60)}} ,
\end{equation}

其中 $\mu_t^{(60)}$ 与 $\sigma_t^{(60)}$ 为价差的滚动 60 日均值与标准差。

\textbf{带宽的理论依据}。论文附录给出：最优无交易带的宽度满足

\begin{equation}
b\;\propto\;\left(\frac{\text{单位交易成本}}{\text{风险厌恶}\times\theta}\right)^{1/3},
\end{equation}

即带宽随成本以三次方根增长\cite{martin2011}。这条关系是本文“反转策略不能频繁穿越阈值”这一工程要求的理论出处。

\textbf{离散化的事件规则}（把连续带宽翻译成可执行的进出场）：

\begin{table}[htbp]
\centering\footnotesize
\caption{期货统计套利的进出场规则}\label{tab:band}
\begin{tabular}{@{}>{\raggedright\arraybackslash}p{2.4cm}>{\raggedright\arraybackslash}p{4.0cm}>{\raggedright\arraybackslash}p{7.6cm}@{}}
\toprule
\textbf{条件} & \textbf{动作} & \textbf{依据} \\
\midrule
$|z_t|>2.0$，无持仓 & 开仓：$z>2$ 空 A 多 B；$z<-2$ 反之 & s-score 进出场范式\cite{avellaneda2010} \\
$|z_t|<0.5$，有持仓 & 平仓 & 同上 \\
$|z_t|>3.5$，有持仓 & 强制止损 & 不收敛防线\cite{leungli2015} \\
持仓 $>3H$ 且未收敛 & 平仓 & 时间止损（资金效率） \\
\bottomrule
\end{tabular}
\end{table}

\textbf{参数边界}。$2.0$ 与 $0.5$ 取自 Avellaneda--Lee 的 s-score 范式（开仓阈值 $2$、平仓阈值 $0.5\sim0.75$）\cite{avellaneda2010}；$3.5$ 的止损阈值为工程标定。按三次方根关系，若实际单边成本从 2 基点升至 6 基点（3 倍），带宽应约放宽至 $3^{1/3}\approx1.44$ 倍，即开仓阈值由 $2.0$ 提高到约 $2.9$。这是本文给出的唯一一条\emph{有理论形式}的参数调整规则，值得在实盘中优先采用。

\subsection{组合与分散}

\textbf{持仓数}：建议同时持有 $N\ge10$ 对。论文指出配对交易的价差交易天然对冲了共同因子暴露，但代价是特质风险无法被市场因子抵消，必须靠配对数量分散（论文执行章）。

\textbf{权重}：按价差波动率倒数归一，使每对的风险贡献大致相等。

\textbf{不做波动率目标化}。与规则二同理：反转类策略不用波动率缩放仓位。论文执行章明确指出对反转策略应把波动率信息用于风险预算（收紧止损、检查冗余），而非信号缩放。

\subsection{失效防线}

\subsubsection*{防线一：不收敛止损}

\begin{itemize}[leftmargin=2.2em]
  \item $|z_t|>3.5$：该对强制平仓；
  \item 持仓超过 $3H$ 仍未回到 $|z|<0.5$：该对平仓；
  \item 单对累计亏损超过该对风险预算的 2 倍：平仓。
\end{itemize}

\subsubsection*{防线二：组合止损与资金冗余}

\begin{itemize}[leftmargin=2.2em]
  \item 组合最大回撤达 $20\%$：敞口减半；
  \item 组合最大回撤达 $30\%$：全部平仓，暂停 40 个交易日；
  \item 预留保证金冗余不低于名义敞口的 $40\%$（期货杠杆下这条尤其重要）。
\end{itemize}

\subsubsection*{防线三：产业链集中度}

同一产业链内的持仓对数不超过总对数的一半。理由是产业链内的品种高度相关，若某一产业链的均衡关系整体失效（如政策冲击），集中持仓会同时受损。

\subsubsection*{共同失效的警告}

论文对偶六的边界条目指出：\textbf{两策略的失效在极端去杠杆事件中会同时发生}——2007 年 8 月中，量化反转基金与动量基金都被冲击，因为冲击来自共同的资金端而非各自的信号端\cite{khandanilo2011,brunnermeier2009}。这意味着若同时运行规则一与规则三，“分散化”在资金冲击下会失效。组合层面必须为此预留额外冗余，而不能假设两个策略互为对冲。

\subsection{成本门槛与容量}

\textbf{成本核算}。单边 2 基点，但必须加上一条修正：\emph{反转类交易的冲击成本高于平均值}，因为交易密集发生在流动性较差的时点。论文审计章把这一点列为已知的未修正偏差。实盘应切换到参与率约束的平方根冲击模型。

\textbf{门槛检查}。若以论文表 5 的 $L=60$ 数据作参照（成本占毛利 $6.3\%$，毛利/成本约 $15.9$），本规则的成本余量在全部门槛检查中最充足。但如前所述，配对版本的实际成本占比需自行回测确定。

\textbf{容量约束}。约束来自两处：其一是成本；其二是产业链内的品种深度——同一产业链内可用于配对的流动性规模有限，且配对交易在两个方向同时占用保证金。

\subsection{参数边界表}

\begin{table}[htbp]
\centering\footnotesize
\caption{规则三参数边界}\label{tab:staparam}
\begin{tabular}{@{}>{\raggedright\arraybackslash}p{2.6cm}>{\raggedright\arraybackslash}p{2.2cm}>{\raggedright\arraybackslash}p{3.2cm}>{\raggedright\arraybackslash}p{6.4cm}@{}}
\toprule
\textbf{参数} & \textbf{取值} & \textbf{证据来源} & \textbf{越界后果} \\
\midrule
信号窗口 & 60（备选 120） & 论文表 5 & 短于 20 毛利夏普已为负；$L=20$ 处于过渡带 \\
形成期 / 交易期 & 60 / 21 交易日 & \cite{gatev2006} 滚动范式 & 重叠则价差平稳性被机械保证（look-ahead） \\
协整显著性 & ADF $<5\%$ 临界值 & \cite{englegranger1987,dickeyfuller1979} & 放宽则无效配对涌入 \\
半衰期区间 & $[5,60]$ 交易日 & 论文附录 A.1 & 下限过低成本激增；上限过高资金效率低 \\
开仓 / 平仓阈值 & $2.0$ / $0.5$ & \cite{avellaneda2010} & 阈值过低频繁交易；过高错过机会 \\
止损阈值 & $3.5$ & 工程标定 & 过宽则单对亏损失控 \\
成本上调的带宽调整 & $b\propto c^{1/3}$ & \cite{martin2011} & 成本涨 3 倍须把阈值提至约 $2.9$ \\
持仓对数 & $\ge10$ & 论文执行章 & 过少则特质风险无法分散 \\
产业链集中度 & $\le50\%$ & 工程判断 & 过度集中则单一产业链冲击致命 \\
保证金冗余 & $\ge40\%$ & \cite{brunnermeier2009} & 不足则可能被强制平仓 \\
\bottomrule
\end{tabular}
\end{table}

\subsection{证据映射表}

\begin{table}[htbp]
\centering\footnotesize
\caption{规则三：决策与论文证据的映射}\label{tab:stamap}
\begin{tabular}{@{}>{\raggedright\arraybackslash}p{3.4cm}>{\raggedright\arraybackslash}p{5.3cm}>{\raggedright\arraybackslash}p{2.6cm}>{\raggedright\arraybackslash}p{2.6cm}@{}}
\toprule
\textbf{决策} & \textbf{论文结论} & \textbf{证据来源} & \textbf{置信度} \\
\midrule
期货截面反转方向有效 & $L=60/120$ 净夏普 $0.72/0.56$、$t=2.59/1.98$ & 论文表 5 & \confmid \\
取 $L=60$ 为主窗口 & 全部门槛检查中成本余量最充足 & 论文表 5 & \confmid \\
弃用 $L\le20$ & 毛利夏普已为负或接近零 & 论文表 5 & \confhigh \\
必须做协整资格检验 & 反转类型是条件性策略 & 论文对偶七 & \confmid \\
半衰期筛选 & OU 离散估计与半衰期定义 & 论文附录 A.1 & \confmid \\
带宽随成本立方根变化 & 最优无交易带宽度关系 & \cite{martin2011} & \confmid \\
在线估计对冲比率 & OU 状态空间配对框架 & \cite{elliott2005} & \confmid \\
s-score 进出场范式 & 因子剥离后的统计套利实现 & \cite{avellaneda2010} & \confmid \\
不用波动率缩放 & 同规则二 & 论文执行章 & \confmid \\
止损与资金冗余 & 不收敛加资金断裂是致命处 & \cite{leungli2015,brunnermeier2009} & \confmid \\
\midrule
配对化的收益水平 & \emph{论文未检验配对版本，不可引用表 5 数字} & — & \conflow \\
开仓阈值 $2.0$ 的具体取值 & \emph{论文未给本文样本下的标定} & — & \conflow \\
对冲比 Kalman 参数 & \emph{论文未检验} & — & \conflow \\
产业链划分方案 & \emph{论文未检验} & — & \conflow \\
\bottomrule
\end{tabular}
\end{table}

表中第一行给 \confmid~而非 \confhigh，反映一个必须正视的事实：论文表 5 支持的是“纯截面排序”形式的反转，而本规则把它改造成了“协整配对”形式。这一步改造有理论依据（对偶七、Avellaneda--Lee 范式），但缺少本文样本下的直接实证。这是本规则与规则二在证据强度上的实质差别。
